\section{Conclusion}

In this paper, we proposed E-PDDL, an Epistemic Planning Domain Definition Language, that represents a unified way to specify \emph{Multi-agent Epistemic Planning} (MEP) problems.
We show that E-PDDL can unify and support leading MEP planners, and provide a parser code to the community that is able to generate input specification for two state of the art epistemic reasoners.
Defining a unified language is the first step towards defining common benchmarks and evaluation metrics for all the MEP community.
Moreover, since the new language follows the well-know language PDDL, E-PDDL may also serve as a gateway for researchers already involved in the area of planing that want to investigate the branch of epistemic reasoning.
Finally defining a unified way of specifying the input for different epistemic solvers can help in designing a general purpose epistemic solver that is automatically able to select the best solving strategy, between existing planners, starting from the problem definition that can now be  ``understood" by every considered approach.
We leave the design of such architecture as future work.

While E-PDDL is able to capture most of the features that can be found in an epistemic domain the language only considers three types of actions: ontic, sensing and announcement.
We plan to expand the language syntax so that is able to accept custom action definitions, specified with a custom event model.
Nonetheless, before introducing such feature we must ensure compatibility with the considered epistemic solvers.
%These task constitute a contribution on its own.
For \rpmep the conversion to explicit \knobel update is handled directly by the parser and, therefore, custom actions definition would only generate different chains of belief formulae.
On the other hand, \EFP relies on an epistemic transition function that expects the actions to be ontic, sensing or announcement and it would not be able to handle custom action types.
That is why before expanding E-PDDL we need to implement in \EFP the possibility of using custom action definitions. 

Finally we plan to refine the parser code to better integrate the concept of partial observability and to be more flexible in the definition of the formulae.
%%%Future works
%disjunctive bf
%Beter integration forall parser
%custom update models
%Integration of common knowlege in formulae